\chapter{W5H Metadata: Where}
\label{chap:W5H_Where}

\section{Purpose}

The \emph{Where} field specifies the spatial, contextual, or domain boundary of
the Engram.

\section{Allowed Contexts}

\begin{itemize}
    \item Physical location (if applicable),
    \item Digital or symbolic domain,
    \item Logical subsystem (e.g., EGO-layer, SEGO-layer),
    \item World-line or narrative context,
    \item Operational domain (task environment or frame).
\end{itemize}

\section{Schema}

\begin{quote}
\textbf{Where Schema}\\
\texttt{where}: 
\begin{itemize}
    \item domain: required (physical / conceptual / symbolic / narrative)
    \item location: optional (text)
    \item frame: optional (cognitive layer or subsystem)
\end{itemize}
\end{quote}

\section{Canonicalization Rules}

\begin{itemize}
    \item The domain must remain consistent with Engram type.
    \item Narrative Engrams must document world-line or branch.
    \item Kernel-bound Engrams may omit ``where'' only if invariant across context.
\end{itemize}
